%%%
%%% v0.1: sent to PCP
%%%

\documentclass[a4paper]{moderncv}
%%%%%%%%%%%%%%%%%%%%%%%%%%%%%%%%%%%%%%%%%%%%%%%%%%%%%%%%%%%%%%%%%%%%%%%
%%% specifying options for CV
%%%%%%%%%%%%%%%%%%%%%%%%%%%%%%%%%%%%%%%%%%%%%%%%%%%%%%%%%%%%%%%%%%%%%%%
\newboolean{cvit}
\newboolean{cvaccademic}
\newboolean{cvfull}
\newboolean{cvpoli}

\setboolean{cvit}{false}
\setboolean{cvaccademic}{false}
\setboolean{cvfull}{false}
\setboolean{cvpoli}{false}

%\extendedversion{La versione estesa del CV \`e disponibile a \url{www.sli-m.com/vernizzi.pdf}}

%%%%%%%%%%%%%%%%%%%%%%%%%%%%%%%%%%%%%%%%%%%%%%%%%%%%%%%%%%%%%%%%%%%%%%%
%%% Misc includes and commands
%%%%%%%%%%%%%%%%%%%%%%%%%%%%%%%%%%%%%%%%%%%%%%%%%%%%%%%%%%%%%%%%%%%%%%%
\usepackage[utf8]{inputenc}
%\usepackage{hyperref}
\usepackage{graphicx}
\usepackage{float}

\moderncvtheme[blue]{classic}

\newcommand{\structure}[1]{\color{sectiontitlecolor}\textbf{#1}\color{black}}

\newcommand{\summary}{Summary}
\newcommand{\education}{Education}
\newcommand{\experience}{Experience}
\newcommand{\teaching}{Teaching}
\newcommand{\research}{Scientific Activity}
\newcommand{\projects}{Research Projects}
\newcommand{\languages}{Languages}
\newcommand{\skills}{Computer Skills}
\newcommand{\interests}{Interests and Activities}
\newcommand{\publications}{Publications (Selection)}
\newcommand{\misc}{Other Information}

\newcommand{\italian}[2]{\ifthenelse{\boolean{cvit}}{#1}{#2}}
\newcommand{\accademic}[1]{\ifthenelse{\boolean{cvaccademic}}{#1}{}}
\newcommand{\work}[1]{\ifthenelse{\boolean{cvaccademic}}{}{#1}}
\newcommand{\full}[1]{\ifthenelse{\boolean{cvfull}}{#1}{}}
\newcommand{\short}[1]{\ifthenelse{\boolean{cvfull}}{}{#1}}
\newcommand{\poli}[2]{\ifthenelse{\boolean{cvpoli}}{#1}{#2}}

\newcommand{\otc}{Open\_TC}

%%%%%%%%%%%%%%%%%%%%%%%%%%%%%%%%%%%%%%%%%%%%%%%%%%%%%%%%%%%%%%%%%%%%%%%
%%% Settings for italian language
%%%%%%%%%%%%%%%%%%%%%%%%%%%%%%%%%%%%%%%%%%%%%%%%%%%%%%%%%%%%%%%%%%%%%%%
\italian{
	\usepackage[italian]{babel}

	\renewcommand{\summary}{Competenze ed interessi}
	\renewcommand{\education}{Formazione}
	\renewcommand{\experience}{Esperienze di lavoro}
	\renewcommand{\teaching}{Attivit\`a didattica}
	\renewcommand{\research}{Attivit\`a scientifica}
	\renewcommand{\projects}{Partecipazione a progetti di ricerca}
	\renewcommand{\languages}{Lingue}
	\renewcommand{\skills}{Competenze informatiche}
	\renewcommand{\interests}{Interessi e attivit\`a}
	\renewcommand{\refname}{Pubblicazioni (Selezione)}
	\renewcommand{\misc}{Altre informazioni}
}
{
	\renewcommand{\summary}{Summary}
	\renewcommand{\education}{Education}
	\renewcommand{\experience}{Work experience}
	\renewcommand{\teaching}{Attivit\`a didattica}
	\renewcommand{\research}{Attivit\`a scientifica}
	\renewcommand{\projects}{Research projects}
	\renewcommand{\languages}{Languages}
	\renewcommand{\skills}{Skills}
	\renewcommand{\interests}{Interessi e attivit\`a}
	\renewcommand{\refname}{Pubblicazioni (Selezione)}
	\renewcommand{\misc}{Languages and other information}
}

%%%%%%%%%%%%%%%%%%%%%%%%%%%%%%%%%%%%%%%%%%%%%%%%%%%%%%%%%%%%%%%%%%%%%%%
%%% Settings for accademic CV
%%%%%%%%%%%%%%%%%%%%%%%%%%%%%%%%%%%%%%%%%%%%%%%%%%%%%%%%%%%%%%%%%%%%%%%
\accademic{
	\usepackage{multibib}
	\italian{
		\newcites{preprint,misc}{{Preprint},{Altro}}
	}
	{
		\newcites{preprint,misc}{{Preprint},{Others}}
	}
}

%%%%%%%%%%%%%%%%%%%%%%%%%%%%%%%%%%%%%%%%%%%%%%%%%%%%%%%%%%%%%%%%%%%%%%%
%%% Settings for work CV
%%%%%%%%%%%%%%%%%%%%%%%%%%%%%%%%%%%%%%%%%%%%%%%%%%%%%%%%%%%%%%%%%%%%%%%
\work{
	\nopagenumbers
}

%%%%%%%%%%%%%%%%%%%%%%%%%%%%%%%%%%%%%%%%%%%%%%%%%%%%%%%%%%%%%%%%%%%%%%%
%%% PERSONAL DATA
%%%%%%%%%%%%%%%%%%%%%%%%%%%%%%%%%%%%%%%%%%%%%%%%%%%%%%%%%%%%%%%%%%%%%%%
\firstname{Davide}
\familyname{Vernizzi}
\title{Curriculum Vit\ae}
\italian{
	\address{Via Belfiore, 42}{10125, Torino, Italia}
}
{
	\address{Via Belfiore, 42}{10125, Torino, Italy}
}
\mobile{+39 329 24 73 484}
\poli {
	\phone{+39 011 2276 810}
}
{ }
\poli {
	\email{davide.vernizzi@polito.it}
}
{
	\email{davide.vernizzi@gmail.com}
}
%\homepage{\url{www.sli-m.com}}

\begin{document}
\maketitle

%%%%%%%%%%%%%%%%%%%%%%%%%%%%%%%%%%%%%%%%%%%%%%%%%%%%%%%%%%%%%%%%%%%%%%%
%%% PUBBLICATIONS
%%%%%%%%%%%%%%%%%%%%%%%%%%%%%%%%%%%%%%%%%%%%%%%%%%%%%%%%%%%%%%%%%%%%%%%
\accademic{
	\nocite{*}
}

\work{
	\nocite{cabiddu10trusted,cesena10anonymous,chiariglione10proposal,avanzi08trace}
}

%%%%%%%%%%%%%%%%%%%%%%%%%%%%%%%%%%%%%%%%%%%%%%%%%%%%%%%%%%%%%%%%%%%%%%%
%%% Competenze ed interessi
%%%%%%%%%%%%%%%%%%%%%%%%%%%%%%%%%%%%%%%%%%%%%%%%%%%%%%%%%%%%%%%%%%%%%%%
\work {
	\vspace{-3.5em}  %%% Hack to fit in 2 pages

	\italian{
	\section{\summary}
        Il mio principale campo di lavoro \`e lo sviluppo delle web applications e
        dei backed e delle API utilizzate da App mobile.
        
        Negli ultimi anni il mio lavoro si \`e svolto principalmente presso 
        Ennova srl, una startup del Politecnico di Torino che sviluppa soluzioni
        per grandi realt\`a quali Telecom Italia, Vodafone, Enel ed Iren Energia.
        Durante la mia collaborazione con Ennova mi sono occupato
        di integrare e consolidare le basi dati dei clienti,
        realizzare servizi a supporto degli IVR di Vodafone
        e curare lo sviluppo di diverse App,
        sia coordinando il lavoro degli sviluppatori,
        sia sviluppando in prima persona le parti piu' critiche delle componenti server.

        Prima di occuparmi di web e mobile application,
        mi sono interessato della sicurezza dei sistemi informativi.
		In questo contesto, ho iniziato ad occuparmi di \emph{Trusted Computing} in
		seguito alla partecipazione al progetto europeo \otc, il cui scopo era
		di sviluppare un framework open source che ne sfruttasse le potenzialit\`a.
		In seguito alla conclusione del progetto \otc, ho iniziato ad occuparmi di \emph{Cloud Computing},
		collaborando al progetto europeo
		TClouds, che aveva come obiettivo lo sviluppo di una piattaforma di Cloud Computing
		resiliente e privacy-enabled.

        Per tutta la durata delle mie collaborazioni, sia al Politecnico di Torino, sia ad Ennova,
        mi sono anche occupato di formare il personale attraverso tesi, tirocini accedemici e formazione
        interna.
	}
	{ %english
        I am currently working as a software engineer at Ennova Srl,
        which was awarded with the prize of Italian Startup of Year 2014.

        I develop web applications, backends, and APIs for mobile applications,
        which are mainly used by telecommunication companies and energy providers,
        such as Telecom Italia, Vodafone, Enel, and Iren Energia.

        While working at Ennova, I have been in charge of integrating clients' databases,
        I developed web services used by clients’ IVR,
        and I coordinated the efforts of Ennova's developers in order to develop various mobile apps.
        Moreover, I personally developed some of the most critical server side building blocks in many projects,
        and I have helped to design Ennova's cloud computing architecture.

        Before working on web and mobile applications, I obtained a Ph.D in security of computer systems.
        In this context, I have participated in the European research project \otc,
        aimed at creating an open-source framework for Trusted Computing.
        Following this, I worked on security of cloud computing, contributing to the research project TClouds,
        whose goal was to create a resilient and privacy-friendly cloud computing platform.

        While holding my positions, both in university and in industry,
        I have always been involved into the training of students or new employees.
	} %english
} %work
\accademic {
	Complete here.
} %accademic

%%%%%%%%%%%%%%%%%%%%%%%%%%%%%%%%%%%%%%%%%%%%%%%%%%%%%%%%%%%%%%%%%%%%%%%
%%% EDUCATION
%%%%%%%%%%%%%%%%%%%%%%%%%%%%%%%%%%%%%%%%%%%%%%%%%%%%%%%%%%%%%%%%%%%%%%%
\section{\education}
\work {
	%%% Dottorato - 15/04/2010
	\italian{
	\cventry{\textbf{15 aprile 2010}}{Dottorato in Ingegneria Informatica e dei Sistemi}{Politecnico di Torino}{}{}{%{XXII Ciclo}{}{%
		Tesi: On Trusted and Privacy-Friendly Network Communications.\newline{}
		Relatore: prof. Antonio Lioy.
	}}{
	\cventry{15 April 2010}{Ph.D. in  Computer Science}{Politecnico di Torino}{}{}{%
		Thesis: On Trusted and Privacy-Friendly Network Communications.\newline{}
		Advisor: prof. Antonio Lioy.
	}}

	%%% Laurea Poli - 04/05/2006
	\italian{
	\cventry{\textbf{5 maggio 2006}}{Laurea Specialistica in Ingegneria Informatica}{Politecnico di Torino}{}
		%{}{
		{\textit{104/110}}{%
		Tesi: Self-adaptive parallel algorithms for computer vision applications.%(Algoritmi adattativi ed applicazioni alla computer vision).
		\newline
		Tesi svolta in collaborazione con l'\'Ecole Nationale Sup\'erieure d'Informatique et des Math\'ematiques Appliqu\'ees de Grenoble (ENSIMAG).\newline{}
		Relatori: prof. Bartrolomeo Montrucchio (POLITO), prof. Jean-Louis Roch (ENSIMAG).
	}}{
	\cventry{5 May 2006}{M.Sc. in Computer Engineering}{Politecnico di Torino}{}{}{%{\textit{104/110}}{%
		Thesis: Self-adaptive parallel algorithms for computer vision applications.
		This thesis was written in collaboration with the \'Ecole Nationale Sup\'erieure d'Informatique et des Math\'ematiques Appliqu\'ees de Grenoble (ENSIMAG).\newline{}
		Advisors: prof. Bartrolomeo Montrucchio (POLITO), prof. Jean-Louis Roch (ENSIMAG).
	}}

	%%% Laurea ENSIMAG - 11/07/2005
	\italian{
	\cventry{\textbf{11 luglio 2005}}{Dipl\^ome d'Ing\'enieur
		(Laurea Specialistica in Ingegneria Informatica)}
		{\newline \'Ecole Nationale Sup\'erieure d'Informatique et des Math\'ematiques Appliqu\'ees de Grenoble (ENSIMAG)}{}{}{%
		Tesi: Self-adaptive parallel algorithms for computer vision applications.\newline
		Relatore: prof. Jean-Louis Roch (ENSIMAG).
	}}{
	\cventry{11 July 2005}{Dipl\^ome d'Ing\'enieur (equivalent to M.Sc. in Computer Engineering)}{\'Ecole Nationale Sup\'erieure d'Informatique et des Math\'ematiques Appliqu\'ees de Grenoble (ENSIMAG)}{}{}{%
		Thesis: Self-adaptive parallel algorithms for computer vision applications.\newline{}
		Advisor: prof. Jean-Louis Roch (ENSIMAG).
	}}

	% Laurea breve - 12/09/2003
	\italian{
	\cventry{\textbf{12 sett. 2003}\hfill}{Laurea in Ingegneria Informatica}{Politecnico di Torino}{}
	%{}{
	{\textit{104/110}}{%
	}}{
	\cventry{12 Sept. 2003}{B.Sc. in Computer Engineering}{Politecnico di Torino}{}{}{%{\textit{104/110}}{%
	}}
} %work
\accademic{
	Complete here.
} %accademic

%%%%%%%%%%%%%%%%%%%%%%%%%%%%%%%%%%%%%%%%%%%%%%%%%%%%%%%%%%%%%%%%%%%%%%%
%%% WORK EXPERIENCE
%%%%%%%%%%%%%%%%%%%%%%%%%%%%%%%%%%%%%%%%%%%%%%%%%%%%%%%%%%%%%%%%%%%%%%%
\section{\experience}

\full{
Le esperienze lavorative seguenti contribuiscono alla mia formazione e dimostrano la mia attitudine ad applicare 
al mondo del lavoro i risultati della mia ricerca, o più in generale le mie competenze scientifiche.\\
} %full

\work {
%        
%		\cventry{\textbf{dal 2010}}{Progetto personale}{sli-m.com, The \structure{sli}de \structure{m}achine \url{http://www.sli-m.com}~}{}{}{%
%			Nel 2010, assieme ad un collega, ho iniziato il progetto \emph{sli-m.com},
%			un servizio web per la produzione e la gestione di presentazioni, basato sull'uso di mind map e \LaTeX/Beamer.
%		}

		% Inizio: 01/11/2011
	\italian {
		\cventry{\textbf{dal 2011}}{Sviluppatore}{Ennova s.r.l, \url{http://www.ennova.it/}}{}{}{%
			Sviluppo di web e mobile application nell'ambito del customer care.
            Integrazione e sviluppo di servizi business e corporate per conto di aziende di telecomunicazioni e provider di energia.
		}
    }
    { %english
		\cventry{\textbf{since 2011}}{Developer}{Ennova s.r.l, \url{http://www.ennova.it/}}{}{}{%
            Development of web and mobile applications for customer care.
            Integration and development of enterprise services on behalf of telephone companies and energy providers.
		}
    }

	\italian {
		% Inizio: 01/07/2010
		\cventry{\textbf{2010 -- 2011}}{Assegnista di ricerca}{Politecnico di Torino, Dip. di Automatica e Informatica}{}{}{%
			In seguito al conseguimento del dottorato, ho continuato la mia attivit\`a di ricerca al
            Politecnico di Torino iniziando ad occuparmi di \emph{Cloud Computing}. 
        }
    }
    { %english
        \cventry{\textbf{2010 -- 2011}}{Postdoctoral researcher}{Politecnico di Torino, Dip. di Automatica e Informatica}{}{}{%
            Following the Ph.D I continued my research activity at Politecnico di Torino, focusing on \emph{Cloud Computing}.
        }
    }

	\italian {
		% Inizio 01/07/2006
		\cventry{\textbf{2006 -- 2010}}{Dottorando}{Politecnico di Torino, Dip. di Automatica e Informatica}{}{}{%
			Attivit\`a di ricerca presso il Computer and Network Security Group.
			L'ambito di ricerca ha riguardato principalmente il \emph{Trusted Computing} e le sue
			applicazioni alle reti comunicazione. 
			\full{
				\`E stato studiato come utilizzare questa tecnologia
				per rendere pi\`u sicure le comunicazioni di rete e ne sono state studiate le problematiche
				collegate alla privacy. Questo lavoro \`e stato svolto principalmente all'interno del progetto
				europe Open\_TC .
			}
			Questa attivit\`a mi ha portato al conseguimento del titolo di Dottore di ricerca.
		}
    }
    { %english
		\cventry{\textbf{2006 -- 2010}}{Research activity}{Politecnico di Torino, Dip. di Automatica e Informatica}{}{}{%
            Research activity at Politecnico di Torino on security of computer systems and network.
            The research activities were mainly focused to \emph{Trusted Computing} and its application
            to network of computers, with a particular attention to privacy.
        }
    }
		
	\italian {
		% Inizio XX/XX/2009
		\cventry{\textbf{2009--2010}}{Collaboratore}{SmartRM S.r.l., \url{http://www.smartrm.com}~}{Torino}{}{%
			Trasferimento tecnologico di 
			parte delle conoscenze maturate durante l'attivit\`a di ricerca. 
			\full{
                In seguito a questa collaborazione \`e stata sviluppata la libreria \emph{Trusted Platform Agent} per lo sviluppo di applicationi
                    che utilizzano \emph{Trusted Computing}.
				I risultati di questo lavoro sono stati pubblicati
				in un articolo scritto per la rivista scientifica \emph{IEEE Software} \cite{cabiddu11trusted}.
				La consulenza, ha portato inoltre ad un contrivuto nella definizione dello 
				standard ISO-MPEG Verify Device Protocol~\cite{chiariglione10proposal}.
			}
		}
    }
    { %english
		\cventry{\textbf{2009--2010}}{Consultant}{SmartRM S.r.l., \url{http://www.smartrm.com}~}{Torino}{}{%
            Consulting activity for SmartRM, a company specialized in developing encryption software
            that allows users to easily protect and share digital content.
            I helped SmartRM to integrate \emph{Trusted Computing} principles to content protection.
        }
    }

	\italian {
		% Inizio 30/06/2009
		\cventry{\textbf{2006}}{Collaboratore}{Politecnico di Torino, Dip. di Automatica e Informatica}{}{}{%
			Sicurezza dei sistemi. Studio circa la possibilit\`a di adottare il protocollo
			\emph{Direct Anonymous Attestation (DAA)} all'interno del protocollo TLS
			effettuato per il progetto europeo \otc.
		}
    }

    \italian {
		% Inizio XX/03/2005
		\cventry{\textbf{2005}}{Stage retribuito}{ID-IMAG, \url{http://bit.ly/id-imag}}{Grenoble, Francia}{}{%
			Studio di nuove tecniche per la parallelizzazione, prove teoriche ed esperimenti pratici dei vantaggi introdotti.
			Applicazione dei risultati a problemi di computer vision.
		}
    }
    { %english
		\cventry{\textbf{2005}}{Internship}{ID-IMAG, \url{http://bit.ly/id-imag}}{Grenoble, Francia}{}{%
            Study and development of parallel algorithms used for computer vision.
        }
    }

		% Inizio XX/06/2004
		% Inizio XX/05/2003
    \italian {
		\cventry{\textbf{2003, 2004}}{Stage}{Dynamic Fun S.r.l., \url{http://www.dynamicfun.com/}~}{Torino}{}{%
			Studio e sviluppo di algoritmi per il riconoscimento facciale.
		}
	} %italian
    { %english
		\cventry{\textbf{2003, 2004}}{Internship}{Dynamic Fun S.r.l., \url{http://www.dynamicfun.com/}~}{Torino}{}{%
            Study and development of algorithms for face detection
        }
    }
} %work
% work


\accademic{
	\italian {
	\section{\research}
		Il mio principale campo di ricerca riguarda la sicurezza dei sistemi informativi.
		In questo contesto, ho iniziato ad occuparmi di \emph{Trusted Computing} in
		seguito alla mia partecipazione al progetto europeo \otc, il cui scopo era
		quello di sviluppare un framework open source che sfruttasse le potenzialit\`a
		offerte da questa tecnologia.

		All'interno del progetto, ho concentrato il mio
		lavoro di ricerca sulle comunicazioni di rete proponendo soluzioni che
		sfruttassero il Trusted Computing per migliorarne la sicurezza. I risultati
		di questo studio sono stati pubblicati in articoli scientifici che illustrano
		come utilizzare il Trusted Computing nelle comunicazioni di rete client/server~\cite{armknecht08efficient},
		alle comunicazioni broadcast~\cite{cesena08towards} e alle reti veicolari~\cite{vernizziXXbolstering}.

		Una parte fondamentale della mia ricerca si \`e anche occupata di studiare gli
		impatti che pu\`o avere il Trusted Computing sulla privacy. Il risultato di
		questa seconda attivit\`a di ricerca \`e stata la proposta di uno schema
		di \emph{autenticazione anonima} basato sui protocolli TLS e \emph{Direct Anonymous Attestation (DAA)},
		pubblicato ad una conferenza scientifica~\cite{cesena10anonymous}.

		% Un riassunto dei contributi principali su questi ambiti \`e disponibile nell'abstract della
		% mia tesi di dottorato, disponibile a \url{xxx}.

		Oltre all'attivit\`a di ricerca, ho contribuito al progetto partecipando
		alla realizzazione dei dimostratori presentati alla Commissione Europea
		durante le revisioni annuali, in particolare ho contribuito in maniera
		assidua allo sviluppo del dimostratore del primo anno.
		La collaborazione \`e poi continuata negli anni successivi partecipando % RIPETIZIONE
		al design complessivo dell'architettura del framework sviluppato.

		In seguito al progetto \otc, il mio lavoro di ricerca si \`e spostato sul cloud
		computing.
	} %italian
	{ %english
		Complete here.
	} %english

	\subsection{\projects}
} %accademic
\work{
	\section{\projects}
} %work

	%%%%%%%%%%%%%%%%%%%%%%%%%%%%%%%%%%%%%%%%%%%%%%%%%%%%%%%%%%%%%%%%%%%%%%%
	%%% Descrione dei vari progetti di ricerca
	%%%%%%%%%%%%%%%%%%%%%%%%%%%%%%%%%%%%%%%%%%%%%%%%%%%%%%%%%%%%%%%%%%%%%%%
\italian {
    \cventry{\textbf{2010-2011}}{TClouds (Trustworthy Clouds)}{}{}{\url{www.tclouds-project.eu} }{%
        Sviluppo di una piattaforma di Cloud Computing resiliente e privacy-enabled.\newline{}
        \full{Progetto finanziato da EC - DG/INFSO - FP7 (ICT-2009-257243).\newline{}}
        \full{Periodo di svolgimento: ottobre 2010 - settembre 2013.}
        \work{Principali partner: IBM, Philips, Ospedale San Raffaele (MI).} 
    }
}
{ %english
    \cventry{\textbf{2010-2011}}{TClouds (Trustworthy Clouds)}{}{}{\url{www.tclouds-project.eu} }{%
        TClouds is EU-funded research project. Its main goal is to develop an advanced cloud infrastructure
        that can deliver computing and storage that achieves a new level of security, privacy, and resilience.
        \newline{}
        \full{EC - DG/INFSO - FP7 (ICT-2009-257243).\newline{}}
        \full{Periodo di svolgimento: ottobre 2010 - settembre 2013.}
        \work{Main partners: IBM, Philips, University of Oxford.} 
    }
}

	\full{
		All'interno di questo progetto, iniziato il 1 ottobre 2010 e da considerarsi come il naturale proseguimento del progetto Open\_TC, 
		mi occupo principalmente del coordinamento del \emph{WP 2.4: Architecture and Integrated Platform}, di cui il nostro gruppo è leader.
		L'obiettivo del WP 2.4 è la definizione dell'architettura e la progettazione del framework che saranno utilizzati per
		combinare i risultati di ricerca di altri workpackage, al fine di realizzare un prototipo integrato di infrastruttura di
		Cloud Computing, su cui verranno dimostrate le applicazioni selezionate all'interno del progetto.\\
		Partecipo inoltre ad altre attivit\`a in altri workpackage.
		Come gruppo di ricerca, ci concentreremo su un sistema di logging e auditing volto a garantire la verificabilit\`a della configurazione del cloud, 
		sia al momento dell'uso delle risorse (attestation), sia in tempi successivi (auditing, forensic), con controllo della disclosure dei dati
		per proteggere la privacy degli attori coinvolti.
		Più in dettaglio mi sto occupando della definizione di
		un sistema di logging privacy-enabled ``by design'' e con funzionalit\`a di controllo d'accessi 
		realizzate con tecniche di broadcast encryption (tema di un precedente lavoro~\cite{cesena08towards}).\\
	} %full

    \italian{
        \cventry{\textbf{2010-2011}}{STORK (Secure Identity Across Borders Linked)}{}{}{\url{www.eid-stork.eu} }{%
            Sviluppo di un sistema europeo interoperabile per l'autenticazione tramite eID.\newline{}
            \full{Progetto finanziato da EC - DG/INFSO - CIP (ICT-PSP-2007-224993).\newline{}}
            \full{Periodo di svolgimento: giugno 2008 - dicembre 2011.}
            \work{Principali partner: 14 paesi membri dell'UE.}
        }
    }
    { %english
        \cventry{\textbf{2010-2011}}{STORK (Secure Identity Across Borders Linked)}{}{}{\url{www.eid-stork.eu} }{%
            STORK project was funded by the European Commission as part of its Competitiveness and Innovation Programme. Its aim is to establish a European electronic ID Interoperability Platform that will allow citizens to use their national eID for authenticating and accessing electronic services across Europe.
            \newline{}
            \work{Main partners: 14 EU Member States.}
        }
    }

	\full{
		Il principale contributo a questo progetto è stata la realizzazione di una demo presentata ad ICT 2010, Bruxelles, Belgio.
		La demo ha la finalit\`a di illustrare ad un cittadino europeo le potenzialit\`a del progetto STORK
		e le nuove possibilit\`a da esso offerte per l'identificazione elettronica,
		ed include un questionario per valutare l'interesse del cittadino verso questa tecnologia.
		I cittadini europei potranno autenticarsi a servizi di nuova generazione,
		come ad esempio l'European Citizen Action Service (ECAS),
		utilizzando le proprie credenziali elettroniche (tipicamente carte d'identit\`a elettroniche).
		Ho inoltre partecipato all'integrazione del servizio ECAS all'interno dell'infrastruttura italiana 
		ed ho sviluppato dei moduli specifici per l'autenticazione di cittadini italiani tramite Carta Nazionale dei Servizi (CNS).
		Questi moduli permettono principalmente di derivare attributi specifici da campi estratti dalle credenziali del cittadino,
		ad esempio permettono di calcolare della data di nascita (e dunque l'et\`a) a partire dal codice fiscale.\\
	} %full

    \italian{
    }
    {
        \cventry{\textbf{2009--2010}}{Trusted Platform Agent}{\url{http://goo.gl/mLodg4} }{}{}{%
		C library (Linux, Windows) for application exploiting Trusted Computing and TPM.
		}
    }

    \italian{
        \cventry{\textbf{2007--2009}}{Open\_TC (Open Trusted Computing)}{}{}{\url{http://goo.gl/IDlKeB} }{%
            Sviluppo di un framework per Trusted Computing, basato su software open source.\newline{}
            \full{Progetto finanziato da EC - DG/INFSO - FP6 (IST-2004-027635).\newline{}}
            \full{Periodo di svolgimento: novembre 2005 - aprile 2009.}
            \work{Principali partner: IBM, HP, Infineon, AMD.}
        }
    }
    {
        \cventry{\textbf{2007--2009}}{Open\_TC (Open Trusted Computing)}{}{}{\url{http://goo.gl/IDlKeB} }{%
            Development of a framework for Trusted Computing based on open source software.
            \newline{}
            \work{Main partners: IBM, HP, Infineon, AMD.}
        }
    }

	\full{
		All'interno del progetto mi sono occupato principalmente dello studio di aspetti
		crittografici legati al Trusted Computing, in primis il protocollo Direct Anonymous Attestation (DAA).
		Il risultato più significativo è stata la definizione di un'estensione al protocollo TLS per
		l'autenticazione anonima attraverso DAA~\cite{cesena10anonymous},
		corredata di implementazione e specifiche~\cite{cesena09ssl,cesena09openssl,cesena09openssla} 
		che ne validano la fattibilit\`a tecnica.
		Un altro risultato importante è stato la realizzazione della Trusted Platform Agent (TPA)~\cite{cabiddu10trusted},
		una libreria open source per il Trusted Computing che semplifica l'interfaccia standardizzata dal 
		Trusted Computing Group ed agevola lo sviluppo di applicazioni che fanno uso di questa tecnologia.
		Ho collaborato infine ad alcune parti di modellazione di componenti del sistema Open\_TC, 
		come ad esempio un sealing proxy~\cite{cesena08secure} per sopperire ad alcune limitazioni
		della funzionalit\`a di sealing (bind crittografico di dati sensibili ad uno stato, o
		configurazione, del sistema) standardizzata dal Trusted Computing Group.\\
	} %full

\accademic{
	\subsection{Partecipazione a Scuole Nazionali e Internazionali COMPLETARE}

	\cventry{2007}{Computer forensics, investigazione penale e criminalit\`a tecnologica. Aspetti informatico-giuridici}
		{Universit\`a degli Studi di Milano}{Milano, Italia}{}
		{Ciclo di seminari organizzato dalla Cattedra di Informatica Giuridica, Facolt\`a di Giurisprudenza.}
	\cventry{XX/200X}{Second European Summer School on Trusted Infrastructure Technologies (ETISS)}{RUB}{Bochum, Germania}{}{}
	\cventry{09/2007}{International School on Foundations of Security Analysis and Design (FOSAD)}{Universit\`a di Bologna}{Bertinoro (FC), Italia}{}{}
	\cventry{09/2007}{International School on Foundations of Security Analysis and Design (FOSAD)}{Universit\`a di Bologna}{Bertinoro (FC), Italia}{}{}

	\subsection{Revisione di Articoli Scientifici COMPLETARE}
	\cvline{Conferenze}{CRISIS 2011 - \url{http://crisis2011.cs.upt.ro/}}
	\cvline{Riviste}{IEEE XXX}

} %accademic


\accademic{
%%%%%%%%%%%%%%%%%%%%%%%%%%%%%%%%%%%%%%%%%%%%%%%%%%%%%%%%%%%%%%%%%%%%%%%
%%% TEACHING EXPERIENCE
%%%%%%%%%%%%%%%%%%%%%%%%%%%%%%%%%%%%%%%%%%%%%%%%%%%%%%%%%%%%%%%%%%%%%%%
\section{\teaching}
	\italian {
		\full{
			A partire dal primo anno del Dottorato di Ricerca ho svolto la seguente attivit\`a didattica.
			Tale attivit\`a si \`e svolta sia presso il Politecnico di Torino (POLITO),
			sia presso aziende esterne per corsi di formazione.\\
		}
		%
		\cventry{\textbf{2009}}{Ingegneria del Software}{POLITO (ING4)}{esercitatore}{}{}
		%
		\cventry{\textbf{2009}}{Sicurezza delle reti ed applicazioni wireless}{POLITO (MASWI)}{docente}{}{}
		%
		\cventry{\textbf{2009}}{Sicurezza delle reti wireless}{POLITO (formazione professionale)}{docente}{}{}
		%
		\work{\cventry{\textbf{2009}}{Trusted Computing}{SmartRM (Torino)}{docente}{}{}}
		%
		\cventry{\textbf{2007, 2008}}{Fondamenti di Informatica}{POLITO (ING3)}{esercitatore}{}{}
		%
		\cventry{\textbf{2007}}{Security of computers systems}{Axis s.r.l. (Torino)}{docente}{}{}

		\section{\languages}
		\cvlanguage{Italiano}{\italian{Lingua madre}{Mother tongue}}{}
		\cvlanguage{English}{}{\italian{Ottima padronanza sia parlato, sia scritto}{Fluent}}
		\cvlanguage{Fran\c{c}ais}{\italian{Ottima conoscenza, sia parlato che scritto}{Fluent}}{}
	} % italian
	{
		\full{
			Starting from my first year as Ph.D. student I held the following teaching
			activities, both as teacher assistant at Politecnico di Torino (POLITO) and
			as teacher to private companies.
		}
		%
		\cventry{\textbf{2009}}{Software engineering}{POLITO (ING4)}{teacher assistant}{}{}
		%
		\cventry{\textbf{2009}}{Wireless network security}{POLITO (MASWI)}{teacher}{}{}
		%
		\cventry{\textbf{2009}}{Wireless network security}{POLITO}{teacher}{}{}
		%
		\work{\cventry{\textbf{2009}}{Trusted Computing}{SmartRM (Torino)}{teacher}{}{}}
		%
		\cventry{\textbf{2007, 2008}}{Introduction to computer science}{POLITO (ING3)}{teacher assistant}{}{}
		%
		\cventry{\textbf{2007}}{Introduction to security of computer systems}{Axis s.r.l. (Torino)}{teacher}{}{}

		\section{\languages}
		\cvlanguage{Italiano}{\italian{Lingua madre}{Mother tongue}}{}
		\cvlanguage{English}{}{\italian{Ottima padronanza sia parlato, sia scritto}{Fluent}}
		\cvlanguage{Fran\c{c}ais}{\italian{Ottima conoscenza, sia parlato che scritto}{Fluent}}{}
	} % english
} % academic

%%%%%%%%%%%%%%%%%%%%%%%%%%%%%%%%%%%%%%%%%%%%%%%%%%%%%%%%%%%%%%%%%%%%%%%
%%% TECH SKILLS
%%%%%%%%%%%%%%%%%%%%%%%%%%%%%%%%%%%%%%%%%%%%%%%%%%%%%%%%%%%%%%%%%%%%%%%
\section{\skills}

\work{
	\italian {
		\cvcomputer{OS}{Linux, Mac OS X, Windows}{Altro}{MySQL, git, UML, \LaTeX, Office}
		\cvcomputer{Programmazione}{PHP, C, Java, Python, Ruby}{Web}{PHP, Javascript, HTML, CSS}
		\cvcomputer{Cloud}{Amazon AWS}{}{}

\full{
		\cventry{Libreria sviluppata}{Trusted Platform Agent}{\url{http://security.polito.it/tc/tpa} }{}{}{%
		Libreria C (Linux, Windows) per applicazioni che usano Trusted Computing e TPM.
		}

		\cventry{Libreria sviluppata}{DAA Toolkit}{\url{http://security.polito.it/tc/daa} }{}{}{%
		Patch ad OpenSSL per autenticazione anonima con TLS e DAA.
		}

		\cventry{Web}{sli-m.com: The \structure{sli}de \structure{m}achine}{\url{http://sli-m.com} }{}{}{%
		Servizio web per la produzione e la gestione di presentazioni, basato su mind map e \LaTeX/Beamer.
		}
}
	} %italian
	{ %english
		\cvcomputer{OS}{Linux, Mac OS X, Windows}{Altro}{MySQL, UML, Vim, \LaTeX, Office}
		\cvcomputer{Programming}{C, Java, Python}{Web}{PHP, Javascript, HTML, CSS}	

		%\cventry{Programming library}{DAA Toolkit}{\url{http://security.polito.it/tc/daa} }{}{}{%
		%Patch for OpenSSL for anonymous authentication using TLS and DAA.
		%}

		%\cventry{Web}{sli-m.com: The \structure{sli}de \structure{m}achine}{\url{http://sli-m.com}}{}{}{%
		%Web service for making presentation basing on mind maps and \LaTeX/Beamer.
		%}
	} %english
} % work


\work{
	%%%%%%%%%%%%%%%%%%%%%%%%%%%%%%%%%%%%%%%%%%%%%%%%%%%%%%%%%%%%%%%%%%%%%%%
	\section{\misc}
	%%%%%%%%%%%%%%%%%%%%%%%%%%%%%%%%%%%%%%%%%%%%%%%%%%%%%%%%%%%%%%%%%%%%%%%
	\italian{
		\cvline{\textbf{Lingue}}{Ottima padronanza di inglese e francese, sia parlato sia scritto.}
		\cvline{\textbf{Didattica}}{Attivit\`a didattica al Politecnico di Torino in qualit\`a di assistente del docente
						e presso aziende private in qualit\`a di docente.}
	} % italian
	{
		\cvline{\textbf{Languages}}{Native language: Italian. Fluent in English and French.}
		\cvline{\textbf{Teaching}}{Teaching at Politecnico di Torino as teacher assistant and teaching to private companies as teacher.}
	} % english
} % work



%%%%%%%%%%%%%%%%%%%%%%%%%%%%%%%%%%%%%%%%%%%%%%%%%%%%%%%%%%%%%%%%%%%%%%%
\makeatletter
\renewcommand*{\bibliographyitemlabel}{\@biblabel{\arabic{enumiv}}}
\makeatother

\accademic{
\section{Produzione Scientifica}
\renewcommand{\refname}{Pubblicazioni}
\bibliographystyle{abbrvurl}
\bibliography{vernizzi}
\nocitepreprint{avanzi08pairing}
\bibliographystylepreprint{abbrvurl}
\bibliographypreprint{vernizzi-extra}
\nocitemisc{chiariglione10proposal,cesena09pairing,elkalam09toward,cesena09ssl,cesena09openssl,cesena09openssla}
\bibliographystylemisc{abbrvurl}
\bibliographymisc{vernizzi-extra}
} %accademic

\accademic{
\bibliographystyle{abbrvurl}
\bibliography{vernizzi,vernizzi-extra}
}

\accademic{
\ \\\ \\
Torino, \today
\flushright
Davide VERNIZZI\hspace*{1cm}
} % accademic

%\vfill
%\hfill Torino, 3/12/2012
%\begin{figure}
%\hfill
%\includegraphics[width=0.4\textwidth]{FirmaPenna}
%\end{figure}


\end{document}

